% part: intuitionistic-logic
% chapter: introduction
% section: bhk-interpretation

\documentclass[../../../include/open-logic-section]{subfiles}

\begin{document}

\olfileid{int}{int}{bhk}

\olsection{બ્રાઉર--હાઇટિંગ--કોલ્મોગોરોવ અર્થઘટન}

\begin{editorial}
  BHK અર્થઘટન વડે સ્ફુરણાવાદી વિધાનોની પ્રામાણ્યતાની
  સાબિતીઓ ગૂંચવણભરી છે; તેમને વધુ સારી રીતે સમજાવવી જોઈએ.
\end{editorial}

સ્ફુરણાવાદી સંયોજકોનું એક અનૌપચારિક રચનાત્મક અર્થઘટન છે, જે
સામાન્ય રીતે બ્રાઉર--હાઇટિંગ--કોલ્મોગોરોવ (BHK) અર્થઘટન તરીકે
ઓળખાય છે. તેમાં ``રચના''નો ખ્યાલ વપરાય છે; તેને રચનાત્મક
સાબિતી સમજી શકાય. (વિધિવત્ !!{derivation} તંત્રની
!!a{derivation} સાથે ગૂંચવણ ન થાય તે માટે BHK અર્થઘટનમાં
``સાબિતી'' શબ્દ વાપરતા નથી.) આ સહજ ખ્યાલને આધારે BHK
અર્થઘટન સ્ફુરણાવાદી સંયોજકોનો અર્થ સમજાવે છે.

\begin{enumerate}
\item પરમાણુક વિધાનની રચના શેને કહેવાય તે આપણે જાણીએ છીએ એવું
  ધારી લઈએ છીએ.
\item $!A_1 \land !A_2$ની રચના એ જોડું $\tuple{M_1, M_2}$ છે,
  જેમાં $M_1$ એ~$!A_1$ની અને $M_2$ એ~$!A_2$ની રચના છે.
\item $!A_1 \lor !A_2$ની રચના એ જોડું $\tuple{s, M}$ છે:
  ક્યાં તો $s$ એ~$1$ છે અને $M$ એ~$!A_1$ની રચના છે, અથવા
  $s$ એ~$2$ છે અને $M$ એ~$!A_2$ની રચના છે.
\item $!A \lif !B$ની રચના એવું વિધેય છે જે~$!A$ની રચનાને
  $!B$ની રચનામાં ફેરવે છે.
\item $\lfalse$ (અસંગતિ)ની કોઈ રચના નથી.
\item $\lnot !A$ને $!A \lif \lfalse$ના સમાનાર્થી તરીકે
  વ્યાખ્યાયિત કરીએ છીએ. એટલે $\lnot !A$ની રચના એવું વિધેય છે
  જે~$!A$ની રચનાને~$\lfalse$ની રચનામાં ફેરવે છે.
\end{enumerate}

\begin{ex}
દાખલા તરીકે $\lnot \lfalse$ લો. તેની રચના એવું વિધેય છે જે
$\lfalse$ની કોઈ પણ રચનાને ઇનપુટ તરીકે લઈને આઉટપુટમાં
ફરી~$\lfalse$ની રચના આપે. સ્પષ્ટપણે, તાદાત્મ્ય
વિધેય~$\Id{}$ આવી રચના છે: $\lfalse$ની રચના~$M$ આપીએ,
તો $\Id{}(M) = M$ ફરી~$\lfalse$ની રચના આપે છે.
\end{ex}

સામાન્ય રીતે, $\lnot !A$નો અર્થ ``$!A$ની રચના અશક્ય છે'' એવો થાય છે.

\begin{ex}
કોઈ પણ વિધાન~$!A$ માટે $!A \lif \lnot \lnot !A$ સાબિત કરીએ;
એ $!A \lif ((!A \lif \lfalse) \lif \lfalse)$ છે. જોઈએ
એવી રચના વિધેય~$f$ છે, જે $!A$ની રચના~$M$ મળતાં
$(!A \lif \lfalse) \lif \lfalse$ની રચના~$f(M)$ આપે.
હવે $f$ એ $(!A \lif \lfalse) \lif \lfalse$ની રચના કેવી
રીતે બનાવે? આપણે એવું વિધેય $g$ વ્યાખ્યાયિત કરવું જોઈએ,
જે $!A \lif \lfalse$ની રચના~$h$ ઇનપુટ તરીકે લે અને
આઉટપુટમાં~$\lfalse$ની રચના આપે. $g$ને આ રીતે વ્યાખ્યાયિત
કરીએ: પહેલાં મળેલી $!A$ની રચના~$M$ પર ઇનપુટ~$h$ લાગુ કરો.
$h$નું આઉટપુટ~$h(M)$ એ $\lfalse$ની રચના છે; તેથી $M$
એ~$!A$ની રચના હોય તો $f(M)(h) = h(M)$ એ~$\lfalse$ની
રચના છે.
\end{ex}

\begin{ex}
$\lnot(!A \land \lnot !A)$, એટલે કે $(!A
\land (!A \lif \lfalse)) \lif \lfalse$, માટે રચના આપીએ.
આ એવું વિધેય~$f$ છે જે $!A \land (!A \lif \lfalse)$ની
રચના~$M$ ઇનપુટ તરીકે લઈને~$\lfalse$ની રચના આપે છે.
સંયોજન $!B_1 \land !B_2$ની રચના એ જોડું $\tuple{N_1, N_2}$
છે, જેમાં $N_1$ એ~$!B_1$ની અને $N_2$ એ~$!B_2$ની રચના છે.
$!B_1 \land !B_2$ની રચનામાંથી અનુક્રમે $!B_1$ અને $!B_2$ની
રચનાઓ પાછી મેળવતાં વિધેયો $p_1$ અને $p_2$ આ રીતે વ્યાખ્યાયિત
કરી શકીએ:
\begin{align*}
  p_1(\tuple{N_1, N_2}) &= N_1\\
  p_2(\tuple{N_1, N_2}) &= N_2
\end{align*}
$f$ આમ કાર્ય કરે છે: પહેલાં તે પોતાના ઇનપુટ~$M$ પર $p_1$
લાગુ કરે છે અને~$!A$ની રચના મેળવે છે. પછી $M$ પર $p_2$
લાગુ કરીને~$!A \lif \lfalse$ની રચના મેળવે છે. આ રચના
પોતે જ વિધેય~$p_2(M)$ છે, જે~$!A$ની રચનાને ઇનપુટ તરીકે
લે તો~$\lfalse$ની રચના આપે છે. બીજા શબ્દોમાં, $p_2(M)$ને
$p_1(M)$ પર લાગુ કરતાં~$\lfalse$ની રચના મળે છે. તેથી
$f(M) = p_2(M)(p_1(M))$ એમ વ્યાખ્યાયિત કરી શકીએ.
\end{ex}

\begin{ex}
$((!A \land !B) \lif !C) \lif (!A \lif (!B \lif !C))$ની
રચના આપીએ. એટલે, $(!A \land !B) \lif !C$ની રચના~$g$ને
$(!A \lif (!B \lif !C))$ની રચનામાં ફેરવતું વિધેય~$f$
આપીએ. રચના~$g$ પોતે $!A \land !B$ની રચનાઓમાંથી~$!C$ની
રચનાઓ આપતું વિધેય છે. \sourcecorrection{OLMD-083}{મૂળમાં
અહીં આઉટપુટ ``$C$ની રચનાઓ'' લખ્યું છે; આસપાસનાં બધાં
સૂત્રોમાં મેટાચલ $!C$ છે, તેથી તે જ રાખ્યું છે.}
$f(g)$નું આઉટપુટ~$h_g$ એવું વિધેય છે જે $!A$ની રચનાઓમાંથી
$!B$ની રચનાઓને~$!C$ની રચનાઓમાં ફેરવતાં વિધેયો આપે છે.

હા, આ ગૂંચવણભર્યું છે. આપણે વિધેય~$h_g$ બનાવવું છે, જે
ઇનપુટ~$g$ માટે $f$નું આઉટપુટ હશે. $h_g$નું ઇનપુટ એ~$!A$ની
રચના~$M$ છે. $h_g(M)$નું આઉટપુટ એવું વિધેય~$k_M$ હોવું જોઈએ
જે~$!B$ની રચના~$N$ને~$!C$ની રચનામાં ફેરવે. હવે
$k_{g,M}(N) = g(\tuple{M, N})$ રાખો. યાદ રાખો કે
$\tuple{M, N}$ એ~$!A \land !B$ની રચના છે. તેથી
$k_{g,M}$ એ $!B \lif !C$ની રચના છે: તે~$!B$ની
રચનાઓ~$N$ને~$!C$ની રચનાઓમાં ફેરવે છે. હવે
$h_g(M) = k_{g,M}$ રાખો. આ વિધેય~$!A$ની રચનાઓ~$M$ને
$!B \lif !C$ની રચનાઓ~$k_{g,M}$માં ફેરવે છે. અંતે
$f(g) = h_g$ રાખો. આ વિધેય $(!A \land !B) \lif !C$ની
રચનાઓ~$g$ને $!A \lif (!B \lif !C)$ની રચનાઓમાં ફેરવે
છે. હાશ!{}
\end{ex}

$!A \lor \lnot !A$ વિધાનને વર્જિત મધ્યનો નિયમ કહે છે.
ચોક્કસ~$!A$ માટે તેને સાબિત કરી શકાય છે (દાખલા તરીકે,
$\lfalse \lor \lnot \lfalse$), પરંતુ સામાન્ય રીતે નહીં.
કારણ કે સ્ફુરણાવાદી વિયોજન માટે બંનેમાંથી કોઈ એક વિયોજ્યની
રચના જરૂરી છે, જ્યારે કેટલાંક વિધાનો હાલ સાબિત પણ કરી શકાતાં
નથી અને ખંડિત પણ કરી શકાતાં નથી (જેમ કે ગોલ્ડબાખની પરિકલ્પના).
તેમ છતાં વર્જિત મધ્યના નિયમને ખંડિત પણ કરી શકાતો નથી:
એટલે $\lnot \lnot (!A \lor \lnot !A)$ સત્ય છે.

\begin{ex}
$\lnot \lnot (!A \lor \lnot !A)$ સાબિત કરવા એવું વિધેય~$f$
જોઈએ જે $\lnot(!A \lor \lnot !A)$ની, એટલે
$(!A \lor (!A \lif \lfalse)) \lif \lfalse$ની રચનાને
~$\lfalse$ની રચનામાં ફેરવે. બીજા શબ્દોમાં, એવું વિધેય~$f$
જોઈએ કે $g$ એ $\lnot(!A \lor \lnot !A)$ની રચના હોય તો
$f(g)$ એ~$\lfalse$ની રચના હોય.

ધારો કે $g$ એ $\lnot(!A \lor \lnot !A)$ની રચના છે; એટલે
એવું વિધેય છે જે $!A \lor \lnot !A$ની રચનાને~$\lfalse$ની
રચનામાં ફેરવે છે. $!A \lor \lnot !A$ની રચના એ જોડું
$\tuple{s, M}$ છે, જેમાં ક્યાં તો $s = 1$ અને $M$
એ~$!A$ની રચના છે, અથવા $s = 2$ અને $M$ એ $\lnot !A$ની
રચના છે. હવે $h_1$ એવું વિધેય હોય જે $!A$ની રચના~$M_1$ને
$!A \lor \lnot !A$ની રચનામાં ફેરવે: તે $M_1$ને
$\tuple{1, M_1}$માં મોકલે છે.
\sourcecorrection{OLMD-084}{મૂળમાં $h_1$નું ઇનપુટ $M_1$
હોવા છતાં આઉટપુટ $\tuple{1, M_2}$ લખાયું છે. વિયોજનની
પ્રથમ શાખાની રચના માટે $\tuple{1, M_1}$ મૂક્યું છે.}
એ જ રીતે $h_2$ એ $\lnot !A$ની રચના~$M_2$ને
$!A \lor \lnot !A$ની રચનામાં ફેરવતું વિધેય હોય: તે
$M_2$ને $\tuple{2, M_2}$માં મોકલે છે.

હવે $k$ને $\comp{h_1}{g}$ માનો: આ વિધેય~$!A$ની રચના
મળે તો~$\lfalse$ની રચના આપે છે; એટલે તે $!A \lif \lfalse$
અથવા $\lnot !A$ની રચના છે. હવે $l$ને $\comp{h_2}{g}$
માનો. આ વિધેયને $\lnot !A$ની રચના આપીએ તો તે~$\lfalse$ની
રચના આપે છે. $k$ એ $\lnot !A$ની રચના હોવાથી, $l(k)$
એ~$\lfalse$ની રચના છે.

આ રીતે $\lnot(!A \lor \lnot !A)$ની રચના~$g$માંથી
~$\lfalse$ની રચના કેવી રીતે મેળવવી તે બતાવ્યું છે.
એટલે $\lnot(!A \lor \lnot !A)$ની રચના~$g$ને
~$\lfalse$ની રચના~$l(k)$માં ફેરવતું વિધેય~$f$ એ
$\lnot\lnot(!A \lor \lnot !A)$ની રચના છે.
\end{ex}

જોઈ શકાય છે કે !!{formula}sની સ્ફુરણાવાદી પ્રામાણ્યતા
બતાવવા BHK અર્થઘટન વાપરવું ઝડપથી ભારરૂપ અને ગૂંચવણભર્યું
બની જાય છે. સદભાગ્યે સ્ફુરણાવાદી તર્ક માટે વધુ સારી
!!{derivation} તંત્રો અને વધુ ચોક્કસ અર્થવિજ્ઞાનલક્ષી
અર્થઘટનો ઉપલબ્ધ છે.
\end{document}
